\section{Conclusion and Future Work}
\label{sec:conclusion}

In this work, we introduced \textbf{GP-BayesMerge}, a mathematically rigorous
Gaussian Process PAC-Bayes framework that addresses the fundamental
vulnerabilities of optimization-based test-time model merging. We exposed the
\emph{Overfitting-Optimizer Paradox}, showing that unconstrained first-order
optimizers easily overfit the transductive noise of small test-time calibration
streams. This overfitting leads to volatile, high-frequency spatial oscillations
across network layers and results in severe generalization collapse on unseen
test data, particularly on SVHN.

By reformulating test-time model merging through the lens of PAC-Bayes
generalization theory, we derived a formal complexity penalty based on the
Kullback-Leibler divergence between the coefficient posterior and a continuous
GP prior over normalized network depth. This formulation naturally simplifies to
a quadratic regularizer governed by the GP precision matrix
$\Sigma_{\ell}^{-1}$, which mathematically unifies proximity constraints
(distance from uniform prior mean) and continuous spatial smoothness
(finite-difference adjacent-layer smoothing). We also established the inverse
depth-scaling rule ($\ell \propto 1/L$) as a critical design guideline for
ultra-deep networks, guaranteeing numerical stability and preserving constant
physical block correlation distances. Our evaluations demonstrate that
GP-BayesMerge completely resolves the Overfitting-Optimizer Paradox, achieving
state-of-the-art accuracy ($84.76\%$) and exceptional stability ($0.37\%$
standard deviation) while outperforming both heuristic spatial smoothing and
rigid subspace constraints.

\textbf{Future Work.} We identify several key directions to extend this
framework. First, while our physical weight-merging validation is evaluated on
pre-trained OpenAI CLIP ViT-B/32 (86M parameters) and CLIP ViT-L/14 (307M parameters)
backbones, a natural extension is to scale GP-BayesMerge to even larger, more
heterogeneous architectures, such as ConvNeXt-XXL, mixture-of-experts (MoE)
models, or multi-modal foundation models, where layers exhibit complex
structural transitions. Second, to scale GP-BayesMerge to ultra-deep
decoder-only Large Language Models (LLMs) like the 32-layer LLaMA-8B or the
80-layer LLaMA-70B, our tridiagonal OU exact inversion provides a powerful
mechanism. We have conducted offline latency benchmarks (reported in Table~\ref{tab:latency} in the Appendix)
showing that our tridiagonal OU formulation scales linearly in $O(L)$ time and
takes less than $0.2$ ms even for $80$ layers, confirming its physical
scalability. In future work, we plan to implement GP-BayesMerge on actual
multi-billion-parameter LLM weight-merging configurations to verify its physical
performance on linguistic and conversational downstream tasks. Third, developing
learned or dynamic non-stationary kernels (such as neural processes or trainable
lengthscales per block) that can automatically discover block boundaries and
adapt spatial correlations based on localized activation statistics.
Specifically, this can be achieved by utilizing localized task activations
during test-time adaptation to compute empirical feature-space distances across
layers. These activation-derived distances can dynamically scale the local GP
lengthscale $\ell_l$ or parameterize a neural covariance kernel, allowing the
model to adaptively partition its layers into functional blocks on-the-fly and
resolve non-stationary architectural transitions. Fourth, resolving the
randomized-to-deterministic PAC-Bayes discrepancy by exploring randomized
posterior evaluation (i.e., sampling merging coefficients $\Lambda \sim Q$ at
inference time) to align empirical execution exactly with the PAC-Bayes
randomized classifier guarantees, which could improve model calibration and
uncertainty estimation under test-time adaptation. Finally, addressing the
fundamental surrogate-to-target risk gap by developing
supervised/self-supervised proxy losses or online test-time calibration
diagnostics that can formally bound target classification error even under
severe, uncalibrated out-of-distribution shifts, as well as extending batch-wise
adaptation to streaming, online Bayesian updates using sequential Monte Carlo or
variational filtering. Crucially, under non-stationary target streams with temporal domain drift, the task-correlation prior $B_{\text{online}}$ can be sequentially updated on-the-fly using sequential Bayesian updates or sliding temporal windows over CKA features, preventing representational lag on dynamically evolving tasks.

Finally, we address the security and potential negative societal impacts of test-time weight interpolation. Since GP-BayesMerge adapts model parameters dynamically based on incoming unlabeled test data, a slow, malicious calibration stream (an adversarial poisoning attack) could attempt to manipulate the learned coefficients to trigger a representational collapse. However, our continuous GP prior $\Sigma_{\ell}^{-1}$ acts as a robust, mathematically principled stabilizer. By heavily penalizing coordinate deviations from the uniform prior mean and enforcing strict spatial smoothness constraints, the GP prior serves as a defensive anchor that constrains malicious coefficient drift, ensuring safe and reliable weight-merging deployments in real-world environments.
